% Part: methods
% Chapter: induction
% Section: relations

\documentclass[../../../include/open-logic-section]{subfiles}

\begin{document}

\olfileid[id]{mth}{ind}{rel}

\olsection{Relasi dan Fungsi}

Ketika kita telah mendefinisikan suatu himpunan objek (seperti bilangan asli
atau suku rapi) secara induktif, kita juga dapat mendefinisikan
\emph{relasi pada} objek-objek tersebut melalui induksi. Misalnya,
perhatikan gagasan berikut: suatu suku rapi~$t_1$ adalah subsuku dari suku
rapi~$t_2$ jika suku pertama muncul sebagai bagian dari suku kedua. Mari kita gunakan
simbol untuk menyatakannya: $t_1 \sqsubseteq t_2$. Tentu saja, setiap suku
rapi merupakan subsuku dari dirinya sendiri: $t \sqsubseteq t$. Kita dapat
memberikan definisi induktif bagi relasi ini sebagai berikut:

\begin{defn}
  Relasi bahwa suatu suku rapi~$t_1$ merupakan subsuku dari~$t_2$,
  $t_1 \sqsubseteq t_2$, didefinisikan melalui induksi pada~$t_2$ sebagai
  berikut:
  \begin{enumerate}
  \item Jika $t_2$ adalah sebuah huruf, maka $t_1 \sqsubseteq t_2$ jika dan
    hanya jika $t_1 = t_2$.
  \item Jika $t_2$ adalah $[s_1 \circ s_2]$, maka $t_1 \sqsubseteq t_2$
    jika dan hanya jika $t_1 = t_2$, $t_1 \sqsubseteq s_1$, atau
    $t_1 \sqsubseteq s_2$.
  \end{enumerate}
\end{defn}

Definisi ini, misalnya, memberi tahu kita bahwa $\mathrm{a}
\sqsubseteq [\mathrm{b} \circ \mathrm{a}]$. Sebab, (2) menyatakan bahwa
$\mathrm{a} \sqsubseteq [\mathrm{b} \circ \mathrm{a}]$ jika dan hanya jika
$\mathrm{a} = [\mathrm{b} \circ \mathrm{a}]$, atau $\mathrm{a}
\sqsubseteq \mathrm{b}$, atau $\mathrm{a} \sqsubseteq \mathrm{a}$. Dua
yang pertama salah: jelas bahwa $\mathrm{a}$ tidak identik dengan
$[\mathrm{b} \circ \mathrm{a}]$, dan menurut~(1), $\mathrm{a}
\sqsubseteq \mathrm{b}$ jika dan hanya jika $\mathrm{a} = \mathrm{b}$,
yang juga salah. Namun, juga menurut~(1), $\mathrm{a} \sqsubseteq
\mathrm{a}$ jika dan hanya jika $\mathrm{a} = \mathrm{a}$, dan ini benar.

Penting untuk diperhatikan bahwa keberhasilan definisi ini bergantung pada
suatu fakta yang belum kita buktikan: setiap suku rapi~$t$ hanya berupa
sebuah huruf, atau terdapat suku rapi $s_1$ dan $s_2$ yang
\emph{ditentukan secara unik} sedemikian sehingga $t = [s_1 \circ s_2]$.
Di sini, ``ditentukan secara unik'' berarti bahwa jika $t = [s_1 \circ
s_2]$, tidak \emph{sekaligus} berlaku $= [r_1 \circ r_2]$ dengan
$s_1 \neq r_1$ atau $s_2 \neq r_2$. Jika hal itu dapat terjadi,
klausa~(2) mungkin bertentangan dengan dirinya sendiri: dengan membaca
$t_2$ sebagai $[s_1 \circ s_2]$, kita mungkin memperoleh $t_1 \sqsubseteq
t_2$, tetapi dengan membaca $t_2$ sebagai $[r_1 \circ r_2]$, kita mungkin
memperoleh bahwa bukan $t_1 \sqsubseteq t_2$. Sebelum membuktikan bahwa hal
ini tidak dapat terjadi, mari kita lihat contoh ketika hal itu
\emph{dapat} terjadi.

\begin{defn}
  Definisikan \emph{suku tanpa kurung} secara induktif sebagai berikut:
  \begin{enumerate}
  \item Setiap huruf merupakan suku tanpa kurung.
  \item Jika $s_1$ dan $s_2$ merupakan suku tanpa kurung, maka
    $s_1 \circ s_2$ merupakan suku tanpa kurung.
  \item Tidak ada yang lain yang merupakan suku tanpa kurung.
  \end{enumerate}
\end{defn}

Contoh suku tanpa kurung ialah $\mathrm{a}$, $\mathrm{b} \circ
\mathrm{d}$, dan $\mathrm{b} \circ \mathrm{a} \circ \mathrm{b}$. Jika
kita mendefinisikan ``subsuku'' bagi suku tanpa kurung seperti di atas,
klausa keduanya akan berbunyi
\begin{center}
  Jika $t_2 = s_1 \circ s_2$, maka $t_1 \sqsubseteq t_2$ jika dan hanya
  jika $t_1 = t_2$, $t_1 \sqsubseteq s_1$, atau $t_1 \sqsubseteq s_2$.
\end{center}

Sekarang $\mathrm{b} \circ \mathrm{a} \circ \mathrm{b}$ berbentuk
$s_1 \circ s_2$ dengan
\begin{align*}
  s_1 & = \mathrm{b} \text{ dan} & s_2 & = \mathrm{a} \circ \mathrm{b}.
\intertext{Untai itu juga berbentuk $r_1 \circ r_2$ dengan}
  r_1 & = \mathrm{b} \circ \mathrm{a} \text{ dan} & r_2 &= \mathrm{b}.
\end{align*}
Apakah $\mathrm{a} \circ \mathrm{b}$ merupakan subsuku dari $\mathrm{b}
\circ \mathrm{a} \circ \mathrm{b}$? Jawabannya ya menurut pembacaan
pertama dan tidak menurut pembacaan kedua.

Sifat bahwa cara suatu suku rapi dibangun dari suku-suku rapi lain bersifat
unik disebut \emph{keterbacaan unik}. Karena definisi induktif bagi relasi
pada objek-objek yang didefinisikan secara induktif semacam ini penting,
kita harus membuktikan bahwa sifat tersebut berlaku.

\begin{prop}
  Andaikan $t$ merupakan suku rapi. Maka $t$ hanya berupa sebuah huruf, atau
  terdapat suku rapi $s_1$, $s_2$ yang ditentukan secara unik sedemikian
  sehingga~$t = [s_1 \circ s_2]$.
\end{prop}

\begin{proof}
  Jika $t$ hanya berupa sebuah huruf, syarat tersebut terpenuhi. Jadi,
  andaikan $t$ bukan hanya berupa sebuah huruf. Dari definisi induktif kita
  dapat mengetahui bahwa $t$ harus berbentuk $[s_1 \circ s_2]$ untuk suatu
  suku rapi $s_1$ dan~$s_2$. Tinggal ditunjukkan bahwa keduanya ditentukan
  secara unik, yakni jika $t = [r_1 \circ r_2]$, maka $s_1 = r_1$ dan
  $s_2 = r_2$.

  Jadi, andaikan $t = [s_1 \circ s_2]$ dan juga $t = [r_1 \circ r_2]$
  bagi suku rapi $s_1$, $s_2$, $r_1$, $r_2$. Kita harus menunjukkan bahwa
  $s_1 = r_1$ dan $s_2 = r_2$. Pertama, $s_1$ dan $r_1$ harus identik;
  sebab jika tidak, yang satu merupakan segmen awal sejati dari yang lain.
  Namun, menurut \olref[sti]{prop:initial}, hal itu mustahil jika $s_1$
  dan~$r_1$ keduanya merupakan suku rapi. Jika $s_1 = r_1$, jelas pula
  bahwa $s_2 = r_2$.
\end{proof}

Kita juga dapat mendefinisikan fungsi secara induktif. Misalnya, kita dapat
mendefinisikan fungsi~$f$ yang memetakan sebarang suku rapi ke kedalaman
maksimum $[\dots]$ yang tersarang di dalamnya sebagai berikut:

\begin{defn}
  \ollabel{defn:depth} \emph{Kedalaman} suatu suku rapi, $f(t)$,
  didefinisikan secara induktif sebagai berikut:
  \[
  f(t) = \begin{cases}
    0 & \text{ jika $t$ adalah sebuah huruf}\\
    \max(f(s_1), f(s_2)) + 1 & \text{ jika $t = [s_1 \circ s_2]$.}
  \end{cases}
  \]
\end{defn}

Misalnya,
\begin{align*}
  f([\mathrm{a} \circ \mathrm{b}]) & =
    \max(f(\mathrm{a}),f(\mathrm{b})) + 1 = \\
  &= \max(0, 0) + 1 = 1, \text{ dan}\\
  f([[\mathrm{a} \circ \mathrm{b}] \circ \mathrm{c}]) & =
    \max(f([\mathrm{a} \circ \mathrm{b}]), f(\mathrm{c})) + 1 = \\
  & = \max(1,0) + 1 = 2.
\end{align*}

Di sini, tentu saja, kita mengasumsikan bahwa $s_1$ dan $s_2$ merupakan
suku rapi dan menggunakan fakta bahwa setiap suku rapi hanya berupa sebuah
huruf atau berbentuk $[s_1 \circ s_2]$. Sekali lagi, penting bahwa suatu
suku dapat memiliki bentuk ini hanya dengan satu cara. Untuk melihat
alasannya, perhatikan kembali suku tanpa kurung yang didefinisikan di atas.
``Definisi'' yang bersesuaian ialah:
\[
  g(t) =
  \begin{cases}
    0 & \text{ jika $t$ adalah sebuah huruf}\\
   \max(g(s_1), g(s_2)) + 1 & \text{ jika $t = s_1 \circ s_2$.}
  \end{cases}
\]
Sekarang perhatikan suku tanpa kurung $\mathrm{a} \circ \mathrm{b} \circ
\mathrm{c} \circ \mathrm{d}$. Suku itu dapat dibaca dengan lebih dari satu
cara, misalnya sebagai $s_1 \circ s_2$ dengan
\begin{align*}
  s_1 & = \mathrm{a} \text{ dan} &
  s_2 & = \mathrm{b} \circ \mathrm{c} \circ \mathrm{d},
\intertext{atau sebagai $r_1 \circ r_2$ dengan}
  r_1 & = \mathrm{a} \circ \mathrm{b} \text{ dan} &
  r_2 &= \mathrm{c} \circ \mathrm{d}.
\end{align*}
Menghitung $g$ menurut pembacaan pertama akan memberikan
\begin{align*}
  g(s_1 \circ s_2) & =
  \max(g(\mathrm{a}), g(\mathrm{b} \circ \mathrm{c} \circ \mathrm{d})) + 1 =\\ & =
  \max(0,2) + 1 = 3
  \intertext{sedangkan menurut pembacaan lainnya kita memperoleh}
  g(r_1 \circ r_2) & =
  \max(g(\mathrm{a} \circ \mathrm{b}), g(\mathrm{c} \circ \mathrm{d})) + 1=\\ &
  = \max(1,1) + 1 = 2.
\end{align*}
Namun, suatu fungsi harus selalu menghasilkan nilai yang unik; jadi,
``definisi'' kita bagi~$g$ sama sekali tidak mendefinisikan suatu fungsi.

\begin{prob}
  Berikan definisi induktif bagi fungsi $l$, dengan $l(t)$ merupakan
  banyaknya simbol dalam suku rapi~$t$.
\end{prob}

\begin{prob}
  Buktikan melalui induksi struktural pada suku rapi $t$ bahwa $f(t) < l(t)$
  (dengan $l(t)$ merupakan banyaknya simbol dalam~$t$ dan $f(t)$ merupakan
  kedalaman $t$ sebagaimana didefinisikan dalam
  \olref[mth][ind][rel]{defn:depth}).
\end{prob}

\end{document}
